% Pista geo-25s-q4b · Geometria
% Procedència: PAU Catalunya 2025 · Convocatòria de setembre · Sèrie 3 · Exercici 4 — Opció B
\documentclass[11pt,a4paper]{article}
\usepackage{pista}
\pistainfo{Catalunya 2025}{Convocatòria de setembre}{Sèrie 3}

\begin{document}

\section*{\textcolor{blauPista}{Exercici 4 --- Opció B}}

\noindent Considereu el pla $\pi: 2x - y + z = 5$ i el punt $P = (0, 1, 3)$.

\subsection*{\textit{a)} \normalfont Comproveu que la distància del punt $P$ al pla $\pi$ és $\dfrac{\sqrt{6}}{2}$. \hspace*{\fill}\textnormal{[0,5~punts]}}

\pista{En primer lloc, escriu el pla en la forma \mbox{$Ax + By + Cz + D = 0$}, passant el $5$ a l'altre costat:
\[
2x - y + z - 5 = 0
\]
A continuació, aplica la fórmula de la distància d'un punt $(x_{0}, y_{0}, z_{0})$ a un pla:
\[
d(P, \pi) = \frac{|Ax_{0} + By_{0} + Cz_{0} + D|}{\sqrt{A^{2} + B^{2} + C^{2}}}
\]
Si et surt una arrel al denominador, racionalitza la fracció: multiplica el numerador i el denominador per aquesta arrel.}

\subsection*{\textit{b)} \normalfont Trobeu l'equació general d'un pla $\pi_{1}$ paral·lel a $\pi$ i que passi pel punt $P$. Quina és la distància entre $\pi_{1}$ i $\pi$? \hspace*{\fill}\textnormal{[0,75~punts]}}

\pista{Aquest apartat et demana dues coses.

\medskip
\textit{El pla $\pi_{1}$.} Com que el pla $\pi_{1}$ és paral·lel a $\pi$, té el mateix vector normal. Per tant, la seva equació és d'aquesta forma:
\[
2x - y + z + D = 0
\]
Troba $D$ fent que el pla passi pel punt $P$.

\medskip
\textit{La distància.} Com que els plans són paral·lels, la distància entre $\pi_{1}$ i $\pi$ és la distància de qualsevol punt de $\pi_{1}$ al pla $\pi$. Quin punt de $\pi_{1}$ coneixes, i quina distància a $\pi$ ja has calculat?}

\subsection*{\textit{c)} \normalfont Trobeu l'equació general d'un segon pla $\pi_{2}$, diferent de $\pi_{1}$, que és paral·lel a $\pi$ i que estigui a una distància $\dfrac{\sqrt{6}}{2}$ de $\pi$. \hspace*{\fill}\textnormal{[1,25~punts]}}

\pista{Fes un dibuix: els plans $\pi_{1}$ i $\pi_{2}$ són a la mateixa distància de $\pi$, un a cada costat. Com a l'apartat \textit{b)}, el pla $\pi_{2}$ té aquesta forma:
\[
2x - y + z + D = 0
\]
A continuació, agafa un punt qualsevol de $\pi$. Per exemple, si fas $x = 0$ i $y = 0$ a l'equació de $\pi$, obtindràs un punt fàcil. Finalment, imposa que la distància d'aquest punt al pla $\pi_{2}$ sigui $\dfrac{\sqrt{6}}{2}$: obtindràs una equació amb un valor absolut. Recorda que una equació com $|X| = 5$ vol dir que $X = 5$ o bé $X = -5$. Una de les dues solucions és el pla $\pi_{1}$ de l'apartat \textit{b)}, i l'altra és el pla $\pi_{2}$ que busques.}

\end{document}
